\capitulo{6}{Discusión}

\section{Discusión de los resultados}

Los resultados obtenidos a través de la web y las posteriores simulaciones demuestran que es posible facilitar la toma de decisiones en el ámbito de la ELA. A continuación, se realiza un análisis de estos hallazgos, interpretando su significado  y contrastándolo con las necesidades del sector.

\subsection{Análisis crítico del feedback de los casos reales}
La validación con los dos pacientes reales cedidos por ELACyL arroja una lectura valiosa sobre cómo los usuarios reciben las recomendaciones de la web en base a su experiencia previa con los SAAC. 

\begin{itemize}
    	\item \textbf{El Paciente PAC-T5NE4D} (caso avanzado) representa al usuario que ya depende por completo de la tecnología. Respondió que ``Sí'' utiliza actualmente los SAAC recomendados (la combinación de lector ocular y programas como \textit{Grid 3} o \textit{Communicator 		5}) y declaró estar ``satisfecho, aunque abierto a informarse sobre otras opciones''. Al ver que la web le sugería sistemas que también encajaban con su perfil (como \textit{Verbo} u \textit{OptiKey}), marcó que ``es posible'' que considere usar otro en el futuro porque encuentra 			interesantes esas alternativas. 
    	Este paciente confirmó que uno de los sistemas ya se le recomendó en el pasado y lo había abandonado porque ``no se ajustaba al entorno''. Esto nos indica que la herramienta aún debe pulirse en lo referido al entorno de uso de los SAAC para evitar recomendaciones inútiles.

    	\item \textbf{El Paciente PAC-7GUYVR} (caso de ELA bulbar) ofrece el punto de vista opuesto. Marcó que ``No'' utiliza actualmente la recomendación (el asistente de voz básico), por lo que su nivel de satisfacción con el sistema se registró como ``No aplica''. Al conservar una buena 		movilidad (manos) y un nivel de autonomía 3, el paciente se sigue comunicando con sus recursos naturales y ve el SAAC como algo secundario. 
    	Sin embargo, que haya puntuado la recomendación con un 6/10 y que valore como ``posible'' cambiar a otro sistema en el futuro demuestra el valor preventivo de la aplicación: aunque ahora mismo no necesite el asistente de voz en su día a día, la web ya le ha dejado abierta la 			puerta a soluciones que le resultan interesantes para cuando la enfermedad avance.
\end{itemize}

Dos pacientes coincidieron en una nota de 6/10. Esto demuestra que la herramienta es realista, no pretende sustituir el criterio de un terapeuta, sino ofrecer una orientación que los propios pacientes y cuidadores validan como adecuada para empezar a explorar opciones.

\subsection{Análisis de los casos de prueba}

\begin{itemize}
    \item El caso de acceso por voz (\textbf{PAC-XPIOTW}) demuestra que el sistema sabe adelantarse a la fatiga en las etapas iniciales. No solo se recomienda un software de recomendación, sino también productos de domótica que dotarán a la persona de una gran autonomía. 
    
    \item El caso de baja tecnología (\textbf{PAC-JOUN1Y}) probó la funcionalidad de la página en casos orientados a personas de avanzada edad, donde el manejo de la tecnología podría suponer una complicación añadida. En el día a día, muchos pacientes sufren fatiga visual ante las pantallas, rechazan la informática (tecnología 0) o necesitan un sistema de emergencia sin baterías por si fallan los de mayor tecnología o se encuentran en un entorno donde no pueden utilizarlo (ducha, lluvia...). 
    
    \item Por último, el caso de validación lingüística (\textbf{PAC-SKNTHB}) demostró una disponibilidad linguistica para que el usuario pueda comunicarse en su lengua materna o aquella con la que se encuentra más cómodo.
\end{itemize}

\subsection{Comparativa con trabajos previos y aportación del proyecto}
Como se inidicaba en el Capítulo 3 - teóricos, la mayoría de los catálogos actuales de SAAC se presentan en formatos estáticos (PDF, guías en papel o webs comerciales). Los profesionales de las asociaciones a menudo deben invertir mucho tiempo cruzando los datos e investigando las opciones.

La principal aportación de este trabajo es romper esa barrera mediante la automatización: se ha conseguido centralizar la información dispersa del mercado en un único punto y cruzarla instantáneamente con el perfil clínico y de accesibilidad del usuario. La herramienta actúa como un puente directo entre la ingeniería del software y el día a día de las asociaciones de pacientes.